\section{R ile çözümlü uygulama}
\label{sec:r-b04}

\textbf{Soru.} $\{2,4,6,8\}$ evreninden geri koymalı iki çekimin bütün
sıralı sonuçlarını oluşturun. Örneklem ortalamasının dağılımını ve standart
hatasını benzetim hatası olmadan hesaplayın.

\begin{lstlisting}[style=prismR]
evren <- c(2, 4, 6, 8)
ciftler <- expand.grid(ilk = evren, ikinci = evren)
ortalamalar <- rowMeans(ciftler)
frekans <- table(ortalamalar)
data.frame(ortalama = as.numeric(names(frekans)),
           frekans = as.integer(frekans),
           olasilik = as.numeric(frekans) / nrow(ciftler))
merkez <- mean(ortalamalar)
varyans <- mean((ortalamalar - merkez)^2)
c(beklenen_deger = merkez, varyans = varyans,
  standart_hata = sqrt(varyans))
sigma <- sqrt(mean((evren - mean(evren))^2))
sigma / sqrt(2)
barplot(prop.table(frekans), xlab = "Orneklem ortalamasi",
        ylab = "Olasilik")
\end{lstlisting}

\begin{outputguidebox}
\textbf{Çözüm ve kontrol değerleri.} Toplam 16 sıralı örneklem vardır.
2--8 arasındaki ortalamaların frekansları $1,2,3,4,3,2,1$'dir.
Dağılımın beklenen değeri 5, varyansı 2,5 ve standart hatası
$\sqrt{2{,}5}=1{,}5811$'dir. Evren standart sapması $\sqrt{5}$ olduğundan
$\sigma/\sqrt{2}$ aynı sonucu verir.
\end{outputguidebox}

\textbf{Yorum.} Bütün olasılık dağılımı listelendiği için varyans 16 böleni
ile hesaplanır. \texttt{var(ortalamalar)} ise 15'e böler ve burada aranan
tam dağılım varyansını vermez. Geri koymasız seçim bu kodla aynı tasarım değildir.

\textbf{Pekiştirme ve yanıt.} $\P(\bar X\geq7)$ değerini
\texttt{mean(ortalamalar >= 7)} ile bulun: $3/16=0{,}1875$.